\section{Guía de generación desde la perspectiva SDE/ODE}

\subsection{Guía en la formulación SDE}

El proceso inverso del modelo difusivo puede formularse equivalentemente como una ecuación diferencial estocástica (SDE) en tiempo continuo. Bajo la formulación SDE, la forma discretizada en el dominio del tiempo del proceso inverso es:

$$x_{t-\Delta t} = x_t + f(x_t, t)\Delta t + g(t)\Delta W$$

donde $f$ contiene la función de señal y $g$ es el coeficiente de ruido. Introducir la guía DPS en la formulación SDE también es directo: después de calcular $x_{t-\Delta t}$ en cada paso, se realiza una actualización basada en el gradiente condicional.

\begin{importantbox}{Equivalencia entre SDE y DDPM}
El paso inverso discreto del DDPM puede considerarse como la discretización en el dominio del tiempo de la SDE continua. Ambas formulaciones son completamente equivalentes: la formulación SDE proporciona un marco de tiempo continuo más natural, mientras que el DDPM es su caso particular en pasos de tiempo discretos. La guía DPS puede aplicarse consistentemente en ambas formulaciones.
\end{importantbox}

\subsection{Viabilidad de la formulación ODE}

Además de la formulación SDE, el proceso inverso del modelo difusivo también puede describirse mediante una ODE de flujo de probabilidad. La forma de la ODE de flujo de probabilidad es:

$$\frac{dx}{dt} = f(x, t) - \frac{1}{2}g(t)^2 \nabla_x \log p_t(x)$$

Su característica principal es la ausencia del término de ruido aleatorio; dada una condición inicial $x_T$, la trayectoria es determinista. Las discusiones en clase indican que la idea de guía de DPS puede aplicarse completamente a la formulación ODE, solo añadiendo actualizaciones de gradiente de guía después de resolver cada paso.

\begin{knowledgebox}{Comparación entre tres formulaciones}
\begin{center}
\begin{tabular}{lccc}
\toprule
Característica & DDPM & SDE & ODE \\
\midrule
Dominio temporal & Discreto & Continuo & Continuo \\
Aleatoriedad & Sí & Sí & No \\
Calidad de muestreo & Alta & Alta & Alta \\
Velocidad de muestreo & Lenta (requiere múltiples pasos) & Lenta & Acelerable \\
Guía DPS & Soportado & Soportado & Soportado \\
Fundamento teórico & Inferencia variacional & Análisis estocástico & Flujo de probabilidad \\
\bottomrule
\end{tabular}
\end{center}
Las tres formulaciones enfrentan el mismo problema práctico en términos de guía: el equilibrio entre la precisión de la aproximación de Tweedie y la intensidad de la guía.
\end{knowledgebox}

\subsection{Equilibrio entre número de pasos temporales e intensidad de guía}

\begin{warningbox}{Más pasos temporales vs. mayor peso de guía}
Durante la inferencia, la guía presenta un equilibrio central crítico:
\begin{itemize}
    \item \textbf{Aumentar el número de pasos temporales}: más pasos implican inyecciones de gradiente de guía más frecuentes, lo que mejora el efecto de guía, pero reduce la velocidad de generación.
    \item \textbf{Aumentar el peso de guía}: reducir los pasos temporales pero aumentar la intensidad de la guía en cada paso provoca que las muestras intermedias se desvíen de la distribución de entrenamiento, haciendo que las predicciones de la red de predicción de ruido sean inexactas y, en última instancia, que la salida resultante se vuelva poco realista.
\end{itemize}
Este es el desafío central de los métodos de guía durante la inferencia: la salida o no cumple con la condición dada, o cumple con la condición pero se vuelve poco realista. Cómo lograr un equilibrio entre realismo y cumplimiento de condiciones es clave para ajustar los parámetros.
\end{warningbox}

Si se desea obtener simultáneamente aceleración (menos pasos) y buenos efectos de guía, podría ser necesario introducir información condicional durante la fase de entrenamiento. Sin embargo, esto perdería la ventaja de la guía en tiempo de inferencia de ``no requerir entrenamiento adicional''.

El ponente señaló que este equilibrio se manifiesta en la práctica de la siguiente manera: cuando el peso de guía es demasiado alto, las muestras intermedias $x_t$ se empujan hacia regiones fuera de la distribución de datos de entrenamiento. En estas regiones fuera de la distribución, la red de predicción de ruido $\epsilon_\theta$ nunca ha visto entradas similares, por lo que sus predicciones se vuelven inexactas. Estas predicciones inexactas se acumularán en los pasos posteriores, lo que lleva a una salida final que, aunque potencialmente cumple con la condición dada, se vuelve visualmente poco realista y genera artefactos.

\begin{knowledgebox}{Sugerencias prácticas de ajuste de parámetros}
Según las discusiones en clase, las siguientes reglas empíricas ayudan a obtener mejores resultados en la práctica:
\begin{itemize}
    \item Utilizar un mayor número de pasos de desruido (por ejemplo, 1000 pasos), permitiendo que la perturbación de cada paso de guía sea menor.
    \item El peso de guía puede variar según el paso temporal: utilizar pesos más bajos en las etapas tempranas (alta ruido) y aumentar gradualmente en las etapas posteriores (bajo ruido).
    \item Intentar simultáneamente ambos samplers, SDE y ODE, para comparar los resultados.
    \item Para la misma condición, realizar múltiples muestreos y seleccionar el mejor (estrategia best-of-N).
\end{itemize}
\end{knowledgebox}

\subsection{Resumen de la sección}

La guía DPS puede aplicarse consistentemente en las tres formulaciones: DDPM, SDE y ODE. En la práctica, es necesario ajustar cuidadosamente el número de pasos temporales y el peso de guía para lograr un equilibrio entre el cumplimiento de la condición y el realismo de la salida.


